% Single-column resume template.
%
% Jinja uses the LaTeX-friendly delimiter set here, so braces and backslashes
% still mean what they mean in LaTeX. Two things to know before editing:
%   - A line starting with two percent signs is a Jinja line statement, so
%     comments in this file use a single percent sign.
%   - The delimiters are live even inside comments, which is why this one
%     describes them instead of showing them.
%
% Injected values are escaped automatically by the environment's finalize hook;
% there is no need to remember an escape filter.

\documentclass[letterpaper,10pt]{article}

\usepackage[top=0.5in,bottom=0.5in,left=0.55in,right=0.55in]{geometry}
\usepackage{titlesec}
\usepackage{enumitem}
\usepackage[hidelinks]{hyperref}
\usepackage{lmodern}

\pagestyle{empty}
\setlength{\parindent}{0pt}
\setlength{\parskip}{0pt}
\raggedbottom
\raggedright

% Section heading: bold, with a rule underneath.
\titleformat{\section}
  {\normalsize\bfseries}
  {}{0em}{}
  [\vspace{-7pt}\rule{\linewidth}{0.6pt}\vspace{-4pt}]
\titlespacing*{\section}{0pt}{9pt}{4pt}

% An entry header: organization on the left, location on the right, then the
% role and dates in italics on the next line.
\newcommand{\entryhead}[2]{\textbf{#1}\hfill #2\par}
\newcommand{\entrysub}[2]{\textit{#1}\hfill \textit{#2}\par\vspace{1pt}}

\newlist{bullets}{itemize}{1}
\setlist[bullets]{
  leftmargin=1.3em,
  labelsep=0.45em,
  itemsep=0.5pt,
  parsep=0pt,
  topsep=2pt,
  partopsep=0pt,
  label=\textbullet
}

\begin{document}

\begin{center}
  {\LARGE\bfseries \VAR{profile.name}}\\[3pt]
  \small
  \BLOCK{if profile.phone}\VAR{profile.phone}\BLOCK{endif}
  \BLOCK{if profile.phone and (profile.email or profile.links)} $|$ \BLOCK{endif}
  \BLOCK{if profile.email}\href{mailto:\VAR{profile.email}}{\VAR{profile.email}}\BLOCK{endif}
  \BLOCK{for link in profile.links}
   $|$ \href{https://\VAR{link}}{\VAR{link}}
  \BLOCK{endfor}
\end{center}

\BLOCK{if education}
\section{Education}
\BLOCK{for school in education}
\entryhead{\VAR{school.school}}{\VAR{school.location}}
\entrysub{\VAR{school.degree}}{\VAR{school.graduation}}
\BLOCK{if school.coursework}
\begin{bullets}
  \item Relevant Coursework: \VAR{school.coursework|join(', ')}
\end{bullets}
\BLOCK{endif}
\BLOCK{endfor}
\BLOCK{endif}

\BLOCK{if skills}
\section{Technical Skills}
\begin{bullets}
\BLOCK{for group in skills}
  \item \textbf{\VAR{group.category}:} \VAR{group['items']|join(', ')}
\BLOCK{endfor}
\end{bullets}
\BLOCK{endif}

\BLOCK{if experience}
\section{Experience}
\BLOCK{for entry in experience}
\entryhead{\VAR{entry.title}}{\VAR{entry.location}}
\entrysub{\VAR{entry.role}}{\VAR{entry.dates}}
\begin{bullets}
\BLOCK{for bullet in entry.bullets}
  \item \VAR{bullet}
\BLOCK{endfor}
\end{bullets}
\BLOCK{endfor}
\BLOCK{endif}

\BLOCK{if projects}
\section{Projects}
\BLOCK{for entry in projects}
\BLOCK{if entry.technologies}
\entryhead{\VAR{entry.title} $|$ \textit{\VAR{entry.technologies}}}{\VAR{entry.dates}}
\BLOCK{else}
\entryhead{\VAR{entry.title}}{\VAR{entry.dates}}
\BLOCK{endif}
\begin{bullets}
\BLOCK{for bullet in entry.bullets}
  \item \VAR{bullet}
\BLOCK{endfor}
\end{bullets}
\BLOCK{endfor}
\BLOCK{endif}

\BLOCK{if leadership}
\section{Leadership \& Involvement}
\begin{bullets}
\BLOCK{for item in leadership}
  \item \VAR{item}
\BLOCK{endfor}
\end{bullets}
\BLOCK{endif}

\end{document}
